\section*{Capítulo \ref{ch:g9:ions} --- Los iones y las disoluciones iónicas}

\begin{solution}{exo:g9:ions:1}
Un \omterm{def:g7:atoms-and-molecules:atom}{átomo} o grupo de \omterm{def:g7:atoms-and-molecules:atom}{átomos} que ha perdido o ganado \omterm{def:g9:inside-the-atom:nucleus}{electrones} y tiene
carga. Un \omterm{def:g9:ions:ion}{catión} es positivo (ha perdido \omterm{def:g9:inside-the-atom:nucleus}{electrones}) y un \omterm{def:g9:ions:ion}{anión},
negativo (ha ganado \omterm{def:g9:inside-the-atom:nucleus}{electrones}).
\end{solution}

\begin{solution}{exo:g9:ions:2}
\ce{Mg^{2+}}: 12 \omterm{def:g9:inside-the-atom:proton}{protones}, 10 \omterm{def:g9:inside-the-atom:nucleus}{electrones}.
\end{solution}

\begin{solution}{exo:g9:ions:3}
\ce{KCl}, \ce{MgCl2}, \ce{CaO}.
\end{solution}

\begin{solution}{exo:g9:ions:4}
El agua salada contiene \omterm{def:g9:ions:ion}{iones}, \ce{Na+(aq)} y \ce{Cl-(aq)}, que se
mueven y transportan carga. Las \omterm{def:g7:atoms-and-molecules:molecule}{moléculas} de azúcar no tienen carga.
\end{solution}

\begin{solution}{exo:g9:ions:5}
\ce{CuSO4}, \ce{AlCl3}, \ce{Na2CO3}.
\end{solution}

\begin{solution}{exo:g9:ions:6}
Hierro(III), \ce{Fe^{3+}}: \ce{Fe^{3+}(aq) + 3OH-(aq) -> Fe(OH)3(s)}.
\end{solution}

\begin{solution}{exo:g9:ions:7}
Sulfato: 2 \omterm{def:g9:inside-the-atom:nucleus}{electrones} de más. Nitrato: 1.
\end{solution}

\begin{solution}{exo:g9:ions:8}
Blanco con nitrato de plata: \ce{Cl-}. Azul con hidróxido de sodio:
\ce{Cu^{2+}}. Cloruro de cobre(II), \ce{CuCl2}.
\end{solution}

\begin{solution}{exo:g9:ions:9}
Cuatro: uno a cada lado (a la izquierda, a la derecha, arriba y abajo).
\end{solution}

\begin{solution}{exo:g9:ions:10}
\ce{Fe^{2+}(aq) + 2OH-(aq) -> Fe(OH)2(s)}.
\end{solution}

\begin{solution}{exo:g9:ions:11}
Cloruro de cinc: los \omterm{def:g9:ions:ion}{iones} sodio no dan \omterm{def:g9:ions:precipitate}{precipitado} con el hidróxido de
sodio, y los \omterm{def:g9:ions:ion}{iones} cinc dan uno blanco:
\ce{Zn^{2+}(aq) + 2OH-(aq) -> Zn(OH)2(s)}.
\end{solution}

\begin{solution}{exo:g9:ions:12}
\ce{CaCl2}: 2000 \omterm{def:g9:ions:ion}{iones} cloruro. Cargas: $1000 \times (+2) + 2000 \times
(-1) = 0$.
\end{solution}

\begin{solution}{pb:g9:ions:1}
\textbf{1.} Hierro(II), \ce{Fe^{2+}}: \ce{Fe(OH)2}.

\textbf{2.} Hierro(III), \ce{Fe^{3+}}: \ce{Fe(OH)3}.

\textbf{3.} El \ce{Fe^{2+}} tiene uno más: $26 - 2 = 24$ \omterm{def:g9:inside-the-atom:nucleus}{electrones},
frente a $26 - 3 = 23$ del \ce{Fe^{3+}}.

\textbf{4.} \ce{Fe^{3+}(aq) + 3OH-(aq) -> Fe(OH)3(s)}.

\textbf{5.} Sale del pozo con \omterm{def:g9:ions:ion}{iones} hierro(II); en contacto con el \omterm{def:g4:air-a-mixture-of-gases:air}{aire}
se convierten en \omterm{def:g9:ions:ion}{iones} hierro(III), que forman el sólido anaranjado.

\textbf{6.} \ce{Fe(OH)3}: $(+3) + 3 \times (-1) = 0$.

\textbf{7.} No en cuanto sale: los \omterm{def:g9:ions:ion}{iones} hierro están disueltos y
atraviesan un filtro. Una vez formado el sólido anaranjado, sí: es un
\omterm{def:g9:ions:precipitate}{precipitado}, que un filtro retiene.

\textbf{8.} No: es una \omterm{def:g6:pure-substances-and-mixtures:mixture}{mezcla} (agua e \omterm{def:g9:ions:ion}{iones} disueltos). Recién sacada y
transparente, es homogénea.

\textbf{9.} $56 \times 1.67 \times 10^{-27} = \qty{9.35e-26}{kg}$.

\textbf{10.} $\qty{0.30}{mg} = \qty{0.30e-6}{kg} = \qty{3.0e-7}{kg}$.

\textbf{11.} Un \omterm{def:g9:ions:ion}{ion} se diferencia del \omterm{def:g7:atoms-and-molecules:atom}{átomo} en dos o tres \omterm{def:g9:inside-the-atom:nucleus}{electrones},
cuya masa es despreciable al lado del \omterm{def:g9:inside-the-atom:nucleus}{núcleo}.

\textbf{12.} $3.0 \times 10^{-7} / 9.35 \times 10^{-26} \approx 3.2
\times 10^{18}$ \omterm{def:g9:ions:ion}{iones} hierro por litro.
\end{solution}
